% !TeX root = ../main.tex

\chapter{群与同态}
\section{子群与正规子群}

% 新概念第一次出现时使用：
% \newterm{唯一键}{中文术语}{English term}
% 后续再次出现时直接写中文，不再重复楷体和英文.

设 \(G\) 是一个群.
我们常常希望从 \(G\) 中取出一部分元素，并要求它们在原来的运算下仍然组成一个群.

\begin{definition}\label{def:subgroup}
若非空子集 \(H\subseteq G\) 在 \(G\) 的运算下本身构成一个群，则称 \(H\) 为 \(G\) 的 \newterm{subgroup}{子群}{subgroup}，记作 \(H\leq G\).
\end{definition}

判断一个非空子集是不是子群时，并不需要逐条检查所有群公理.

\begin{lemma}\label{lem:subgroup-test}
设 \(H\) 是 \(G\) 的非空子集.
若对任意 \(x,y\in H\) 都有 \(xy^{-1}\in H\)，则 \(H\leq G\).
\end{lemma}

\begin{proof}
取 \(x\in H\).
令 \(y=x\)，则 \(e=xx^{-1}\in H\).
再令 \(x=e\)，可得 \(y^{-1}\in H\).
最后，对任意 \(x,y\in H\)，由于 \(y^{-1}\in H\)，应用假设得到 \(x(y^{-1})^{-1}=xy\in H\).
因此 \(H\) 在乘法与取逆下封闭，从而是 \(G\) 的子群.
\end{proof}

有些子群与整个群的乘法结构特别相容，这些子群正是构造商群所需要的对象.

\begin{definition}\label{def:normal-subgroup}
若子群 \(N\leq G\) 满足对任意 \(g\in G\) 都有 \(gNg^{-1}=N\)，则称 \(N\) 为 \(G\) 的 \newterm{normal-subgroup}{正规子群}{normal subgroup}，记作 \(N\trianglelefteq G\).
\end{definition}

\begin{example}\label{ex:abelian-normal}
若 \(G\) 是交换群，则 \(G\) 的每个子群都是正规子群.
事实上，对任意 \(g\in G\) 与 \(n\in N\)，都有 \(gng^{-1}=n\).
\end{example}

\begin{remark}\label{rem:normal}
子群是否正规依赖于它嵌入整个群的方式.
同一个抽象群可以在某个群中作为正规子群出现，也可以在另一个群中作为非正规子群出现.
\end{remark}

% 习题直接放在本小节末尾，不单独建立“习题”小节.
\begin{exercises}
    \item 证明任意群 \(G\) 的中心 \(Z(G)\) 都是正规子群.
    \item 设 \(H,K\leq G\)，证明 \(H\cap K\leq G\).
\end{exercises}

% ================================================================
%                     商群与群同态
% ================================================================
\section{商群与群同态}

正规子群允许我们把同一个陪集中的元素视为同一个对象.

\begin{definition}\label{def:quotient-group}
若 \(N\trianglelefteq G\)，则所有陪集
\begin{align*}
    G/N=\{gN\mid g\in G\}
\end{align*}
在乘法 \((gN)(hN)=ghN\) 下构成的群称为 \newterm{quotient-group}{商群}{quotient group}.
\end{definition}

接下来考虑两个群之间保持运算的映射.

\begin{definition}\label{def:homomorphism}
设 \(G\) 与 \(H\) 是群.
若映射 \(\varphi:G\to H\) 满足
\begin{align*}
    \varphi(xy)=\varphi(x)\varphi(y)
\end{align*}
对所有 \(x,y\in G\) 都成立，则称 \(\varphi\) 为 \newterm{homomorphism}{群同态}{group homomorphism}.
\end{definition}

与一个群同态相伴的两个基本对象，是被送到单位元的元素和真正被映到的元素.

\begin{definition}\label{def:kernel-image}
群同态 \(\varphi:G\to H\) 的 \newterm{kernel}{核}{kernel} 与 \newterm{image}{像}{image} 分别定义为
\begin{align}
    \ker\varphi&=\{g\in G\mid \varphi(g)=e_H\},
    &
    \operatorname{im}\varphi&=\{\varphi(g)\mid g\in G\}.
    \label{eq:kernel-image}
\end{align}
\end{definition}

核不仅记录同态把哪些元素压成了单位元（见式~\eqref{eq:kernel-image}），它还精确描述了同态在定义域中丢失的信息.

\begin{proposition}\label{prop:kernel-normal}
对任意群同态 \(\varphi:G\to H\)，都有 \(\ker\varphi\trianglelefteq G\).
\end{proposition}

\begin{proof}
取 \(x\in\ker\varphi\) 与 \(g\in G\).
则
\begin{align*}
    \varphi(gxg^{-1})
    =\varphi(g)\varphi(x)\varphi(g)^{-1}
    =e_H.
\end{align*}
因此 \(gxg^{-1}\in\ker\varphi\)，所以 \(\ker\varphi\) 是正规子群.
\end{proof}

\begin{theorem}\label{thm:first-iso}
设 \(\varphi:G\to H\) 是群同态.
则存在自然同构
\begin{align*}
    G/\ker\varphi\cong\operatorname{im}\varphi.
\end{align*}
\end{theorem}

这个结论通常称为群的第一同构定理，关于这一基本结果也可参见 \cite{lang}.
它说明任何群同态都可以分解为先取商，再把所得商群同构到像上.
完整证明放在本小节末、习题之前，以展示较长证明的双向跳转方式.

\longprooflink{first-iso-proof}{查看完整证明}

\begin{corollary}\label{cor:injective-kernel}
群同态 \(\varphi:G\to H\) 是单射，当且仅当 \(\ker\varphi=\{e_G\}\).
\end{corollary}

\begin{proof}
由 \autoref{thm:first-iso}，商群 \(G/\ker\varphi\) 与 \(\operatorname{im}\varphi\) 同构.
因此 \(\varphi\) 不识别任何两个不同元素，当且仅当核中只有单位元.
\end{proof}

% ================================================================
%                        交换图
% ================================================================
第一同构定理也可以用交换图表示：
\[
\begin{tikzcd}
G \arrow[r,"\varphi"] \arrow[d,"\pi"'] &
\operatorname{im}\varphi \\
G/\ker\varphi \arrow[ur,"\overline{\varphi}"',"\sim"]
\end{tikzcd}
\]
其中 \(\pi\) 是自然投影，而 \(\overline{\varphi}\) 是定理给出的同构.

% ================================================================
%                         例与问题
% ================================================================
\begin{example}\label{ex:modn}
考虑群同态
\begin{align*}
    \varphi:\mathbb Z&\longrightarrow\mathbb Z/n\mathbb Z,\\
    m&\longmapsto\overline m.
\end{align*}
它的核是 \(n\mathbb Z\)，像是整个 \(\mathbb Z/n\mathbb Z\).
因此 \autoref{thm:first-iso} 给出
\begin{align*}
    \mathbb Z/n\mathbb Z\cong \operatorname{im}\varphi.
\end{align*}
\end{example}

\begin{problem}\label{prob:sign}
设
\begin{align*}
    \operatorname{sgn}:S_n\longrightarrow\{\pm1\}
\end{align*}
为符号同态.
求它的核，并利用 \autoref{thm:first-iso} 描述对应的商群.
\end{problem}


% 较长证明放在本小节末、习题之前.
\begin{longproof}{first-iso-proof}{\autoref{thm:first-iso}}
由 \autoref{prop:kernel-normal}，\(\ker\varphi\trianglelefteq G\)，所以商群 \(G/\ker\varphi\) 有定义.
定义
\begin{align*}
    \overline{\varphi}:G/\ker\varphi&\longrightarrow\operatorname{im}\varphi,\\
    g\ker\varphi&\longmapsto\varphi(g).
\end{align*}

首先验证这个映射定义良好.
若
\begin{align*}
    g\ker\varphi=h\ker\varphi,
\end{align*}
则 \(h^{-1}g\in\ker\varphi\)，从而
\begin{align*}
    \varphi(h)^{-1}\varphi(g)
    =\varphi(h^{-1}g)
    =e_H.
\end{align*}
因此 \(\varphi(g)=\varphi(h)\)，所以 \(\overline{\varphi}\) 与陪集代表元的选择无关.

其次，\(\overline{\varphi}\) 是群同态，因为
\begin{align*}
    \overline{\varphi}\bigl((g\ker\varphi)(h\ker\varphi)\bigr)
    =\varphi(gh)
    =\varphi(g)\varphi(h).
\end{align*}

它显然是满射，因为 \(\operatorname{im}\varphi\) 中的每个元素都具有形式 \(\varphi(g)\).

最后验证单射.
若
\begin{align*}
    \overline{\varphi}(g\ker\varphi)=e_H,
\end{align*}
则 \(\varphi(g)=e_H\)，故 \(g\in\ker\varphi\).
于是 \(g\ker\varphi=\ker\varphi\)，所以 \(\overline{\varphi}\) 的核只有单位陪集.

因此 \(\overline{\varphi}\) 是同构，从而
\begin{align*}
    G/\ker\varphi\cong\operatorname{im}\varphi.
\end{align*}
\end{longproof}

% 本小节的习题紧接在正文与长证明之后.
\begin{exercises}
    \item 设 \(\varphi:G\to H\) 为群同态，证明 \(\operatorname{im}\varphi\leq H\).
    \item 设 \(N\trianglelefteq G\)，证明自然投影 \(G\to G/N\) 的核恰好是 \(N\).
    \item 设 \(\operatorname{sgn}:S_n\to\{\pm1\}\) 为符号同态，求它的核.
\end{exercises}
